\section{Related work: a classified review of 20 studies}
\label{sec:literature}
\subsection{Review framework and classification}
We group 20 studies on prompt optimization, evaluation and evolutionary selection
into five categories according to the decision they primarily address: candidate
generation and refinement (7), cost and evaluation allocation (5), structured
search and operator control (2), evaluation robustness (3), and noisy evolutionary
selection (3). The three surveys provide broader context
\cite{survey_li,survey_cui,survey_ramnath}; NSGA-II is foundational algorithmic
background~\cite{nsga2}. They are not counted among the 20 studies in
Table~\ref{tab:twenty-studies}.

Some methods span more than one category. Table~\ref{tab:twenty-studies}
groups them by their main decision, then compares their mechanisms with PACE.
The final column gives our interpretation of that connection. Conference and
journal records are distinguished from preprints in the references. Because
benchmarks, models, operators and accounting conventions differ, no cross-paper
accuracy ranking is constructed.

\subsection{Candidate generation, refinement and structured search}
\paragraph{Evolutionary and self-referential operators.}
EvoPrompt connects LLM-generated language variation with a population-based
evolutionary search over discrete instructions~\cite{evoprompt}. Promptbreeder
adds a second adaptive level: mutation prompts themselves evolve alongside task
prompts~\cite{promptbreeder}. Their common lesson for this study is that survival
has consequences beyond the current ranking, because selected instructions shape
future candidates. Their operator designs are substantially richer than PACE's
rule-based mutations. PACE uses simpler rule-based mutations to keep the survivor comparison
controlled. Results from these different operator families cannot be compared
directly.

\paragraph{Optimization histories and latent search.}
OPRO supplies an LLM with previous candidate solutions and their values, allowing
that history to guide the next proposals~\cite{opro}. InstructZero instead
optimizes a low-dimensional soft prompt applied to an open model, which generates
discrete instructions for a black-box target; Bayesian optimization uses the
target's evaluations~\cite{instructzero}. These approaches differ in search
representation and update mechanism. Both illustrate that a prompt optimizer's
information state matters. PACE's contribution lies in filtering removal decisions
within a population, rather than replacing the proposal mechanism with either
history-conditioned generation or a continuous surrogate search.

\paragraph{Critique, residual structure, and Pareto reuse.}
PromptWizard iteratively critiques and synthesizes instructions and in-context
examples~\cite{promptwizard}. RiOT generates diverse refinements from text
gradients and uses residual connections to retain useful content while limiting
semantic drift~\cite{riot}. GEPA reflects on system trajectories and combines
complementary lessons from a Pareto frontier of attempts~\cite{gepa}. These methods
make feedback more informative than a single scalar score. However, useful
feedback and uncertain numerical rankings address different bottlenecks. Our
research question is whether removal should wait for paired evidence after the
chosen evaluations have been observed. The present experiment addresses this survival decision while keeping the
feedback and proposal mechanism fixed.

\paragraph{Structured search and operator control.}
AutoPDL casts prompting patterns and demonstrations as a structured AutoML space
and uses successive halving to allocate effort~\cite{autopdl}. Its human-readable
program outputs show that the optimization object can be more complex than a
single instruction string. The evolutionary toolbox of Grie{\ss}haber et al.
separates evolutionary stages and studies operator verification, human feedback
and more efficient evaluation~\cite{toolbox}. Both make process decomposition
useful for diagnosis. PACE adopts that analytical separation between generation,
evaluation and survival, while empirically changing only its survivor rule
within the frozen implementation.

\subsection{Cost objectives and evaluation allocation}
CAPO is the closest architectural motivation: it combines evolutionary prompt
optimization, racing and prompt-length considerations~\cite{capo}. MO-CAPO makes
the accuracy--inference-cost frontier explicit and studies budget allocation
within multi-objective search~\cite{mocapo}. They motivate evaluating both the cost
of the search process and the cost of the deployed solution. PACE follows that
distinction but uses a small fixed protocol and holds its evaluation policy
constant in the primary comparison. A positive accuracy difference accompanied
by higher token cost is consequently a trade-off, not unconditional dominance.

BudgetAPO directly targets small evaluation budgets using noise-adaptive slices,
paired accept/reject comparisons and reflective rewriting~\cite{budgetapo}. Pairing is already part of BudgetAPO, so it provides a direct precedent
for paired comparisons under tight budgets. The
specific distinction is a gate for population survival coupled to an explicit
bounded archive and separate capacity-eviction logs. BudgetAPO is cited as a preprint. Its different proposal and evaluation
policies would need to be controlled in a direct empirical comparison.

IPOMP selects representative evaluation examples through semantic clustering,
boundary analysis and performance-guided refinement~\cite{ipomp}. Its focus is
which data support an informative comparison. That consideration matters for
interpreting bootstrap evidence: a narrow or unrepresentative observation set
can give a stable-looking ranking that transfers poorly. PACE uses randomly
shared samples and does not implement IPOMP's data selection. Combining adaptive
example selection with the gate would require reconsidering what distribution
the paired statistic estimates.

IAPO jointly considers prompt choice and inference scale under user preferences
and an inference budget, with a sequential-trimming algorithm and finite-budget
error analysis~\cite{iapo}. This extends the decision space beyond the fixed
greedy-decoding configuration used here. Its guarantees apply to its own
formulation and algorithm; they do not transfer to PACE's percentile bootstrap.
For PACE, both the deployment preference and greedy decoding configuration
are fixed before winner selection.

\subsection{Evaluation robustness and noisy selection}
\paragraph{Model dependence and evaluation design.}
Revisiting OPRO reports limitations when relatively small LLMs act as optimizers
and recommends accounting for model ability and computational
cost~\cite{opro_revisit}. It cautions against treating a successful optimization
strategy on one model as universally effective. Our experiment uses one 7B model and rule-based proposals. It leaves open how
the survivor rule would behave with larger models or stronger proposal operators.

Yang et al. provide a unified interpretation of probability-based prompt
selection through mutual information and study calibration by
marginalization~\cite{probsel}. Their score and supervision assumptions differ
from PACE's observed binary correctness. The scoring rule, calibration and access to model outputs are part of
the selection problem. PACE does not assume token-probability access, and no
probability-based baseline is evaluated here.

Optimization before Evaluation shows that optimizing prompts separately for
models can alter model rankings~\cite{evalbefore}. A common unoptimized prompt
is therefore not always a neutral comparison. In this report, both methods use
the same target model and task, and each final winner is selected using search
observations only. This controls one comparison axis, while the shared held-out
set and narrow operator family still limit external conclusions.

\paragraph{Different sources of evolutionary selection noise.}
Depth over Fidelity studies the tension between denoising rankings and retaining
more evolutionary updates under a fixed budget. Its probabilistic elite
membership replaces hard rank weights with conditional expected weights and is
instantiated through residual bootstrapping~\cite{depth}. This is a close
conceptual alternative to delaying hard elimination. It concerns evolution
strategies in its own search spaces, whereas PACE uses discrete prompt strings,
a bounded archive and deterministic truncation after its gate.

Gulayeva et al. analyze \emph{selection noise} as genetic drift caused by stochastic
reproduction, including takeover-time and reproduction-rate measures
\cite{selectionnoise}. That is different from noisy fitness estimation from a
small example sample. Protecting a prompt from uncertain removal leaves stochastic parent sampling
and its effect on diversity unchanged.

Martin and Collins allocate sampling time according to predicted sorting
difficulty in noisy robot optimization with adaptive-sampling
CMA-ES~\cite{robotnoise}. Their work illustrates that evaluation precision can be
adapted to the selection problem. Its observations, continuous policy space and
physical time costs differ from prompt--article queries. We use it as motivation
for the depth--precision trade-off, not as a directly comparable benchmark.

\subsection{Synthesis and the precise research gap}
The review identifies several established ingredients: language-based variation,
reflection, cost-aware multi-objective search, paired comparisons, adaptive
allocation and uncertainty-aware selection. PACE combines these ingredients
to test how a
particular paired-evidence survival rule interacts with finite retention capacity
when the surrounding prompt-search implementation is held fixed.

The experiment therefore records gate eliminations and capacity evictions
separately, matches total retention capacity online, and reports retention
alongside held-out performance. These measurements test whether the extra
protection changes the outcome of search. The code and per-seed records allow
each comparison to be checked.

\begingroup\small
\setlength{\LTleft}{0pt}\setlength{\LTright}{0pt plus 1fill}
\begin{longtable}{@{}>{\raggedright\arraybackslash}p{2.3cm}>{\raggedright\arraybackslash}p{4.1cm}>{\raggedright\arraybackslash}p{6.5cm}@{}}
\caption{Classification and analytical comparison of 20 core studies. Category
counts are 7/5/2/3/3. The PACE connection is a qualitative synthesis; these studies
are not empirical baselines in the present experiment.}\label{tab:twenty-studies}\\
\toprule Study & Main mechanism or decision & Connection and boundary for PACE\\\midrule
\endfirsthead
\multicolumn{3}{l}{\textit{Table~\thetable\ continued: classification of 20 core studies}}\\
\toprule Study & Main mechanism or decision & Connection and boundary for PACE\\\midrule
\endhead
\midrule\multicolumn{3}{r}{\textit{Continued on next page}}\\\endfoot
\bottomrule\endlastfoot
\multicolumn{3}{@{}l}{\textbf{A. Candidate generation and refinement (7)}}\\*[3pt]
EvoPrompt (2024)~\cite{evoprompt} & LLM variation inside evolutionary prompt search. & Establishes population-based discrete search; PACE studies the later survival decision.\\[3pt]
Promptbreeder (2024)~\cite{promptbreeder} & Co-evolution of task prompts and mutation prompts. & Shows that retained prompts affect future variation; PACE does not evolve mutation instructions.\\[3pt]
OPRO (2024)~\cite{opro} & History-conditioned LLM proposals from scored solutions. & Uses optimization feedback; differs from a population-level removal gate.\\[3pt]
InstructZero (2024)~\cite{instructzero} & Bayesian optimization of a soft prompt that generates discrete instructions. & Black-box target optimization with a different representation and proposal mechanism.\\[3pt]
PromptWizard (2025)~\cite{promptwizard} & Critique and synthesis of instructions and examples. & Richer feedback and variation; no head-to-head inference from PACE's small search.\\[3pt]
RiOT (2025)~\cite{riot} & Text-gradient refinement, residual connections and a search tree. & Addresses diversity and semantic drift rather than the same bounded survivor rule.\\[3pt]
GEPA (2026)~\cite{gepa} & Reflection on trajectories and Pareto reuse of lessons. & Reuses complementary successes; PACE uses binary paired observations and a simpler operator.\\[3pt]
\multicolumn{3}{@{}l}{\textbf{B. Cost objectives and evaluation allocation (5)}}\\*[3pt]
CAPO (2025)~\cite{capo} & Evolutionary search, racing and length-aware optimization. & Closest architectural motivation; PACE isolates a survivor-rule contrast.\\[3pt]
MO-CAPO (2026)~\cite{mocapo} & Multi-objective performance--cost search with budget allocation. & Motivates deployment trade-offs; PACE does not claim a better frontier across tasks.\\[3pt]
BudgetAPO (2026)~\cite{budgetapo} & Noise-adaptive evaluation slices and paired accept/reject decisions. & Direct precedent for pairing under tight budgets; PACE adds its particular archive/survival integration.\\[3pt]
IPOMP (2025)~\cite{ipomp} & Representative data selection refined by model performance. & Changes which examples inform comparisons; not implemented in PACE's random shared sampling.\\[3pt]
IAPO (2026)~\cite{iapo} & Joint prompt and inference-scale selection; sequential trimming. & Makes budget and preference explicit; its error guarantees are not PACE guarantees.\\[3pt]
\multicolumn{3}{@{}l}{\textbf{C. Structured search and operator control (2)}}\\*[3pt]
AutoPDL (2025)~\cite{autopdl} & Structured prompting-program search with successive halving. & Broadens the search object beyond strings; highlights allocation as a separate design choice.\\[3pt]
Evolutionary toolbox (2025)~\cite{toolbox} & Stage decomposition, operator checking, feedback and evaluation improvements. & Supports modular analysis; PACE freezes these surrounding components.\\[3pt]
\multicolumn{3}{@{}l}{\textbf{D. Evaluation robustness and model dependence (3)}}\\*[3pt]
Revisiting OPRO (2024)~\cite{opro_revisit} & Tests optimizer limitations in smaller LLMs. & Motivates model-specific conclusions and simple baselines.\\[3pt]
Probability-based selection (2024)~\cite{probsel} & Unified mutual-information view and calibration. & Selection depends on scoring assumptions; PACE uses correctness rather than token probabilities.\\[3pt]
Optimization before Evaluation (2025)~\cite{evalbefore} & Examines ranking changes after model-specific prompt optimization. & Motivates controlled evaluation; does not resolve PACE's shared-test-set uncertainty.\\[3pt]
\multicolumn{3}{@{}l}{\textbf{E. Noisy evolutionary selection (3)}}\\*[3pt]
Depth over Fidelity (2026)~\cite{depth} & Probabilistic elite membership under fixed evaluation budgets. & Integrates rank uncertainty through weights, unlike PACE's gate and archive.\\[3pt]
Selection-noise analysis (2025)~\cite{selectionnoise} & Genetic drift in stochastic selection schemes. & Reproduction noise differs from noisy estimated fitness; PACE addresses only part of the problem.\\[3pt]
Adaptive-sampling CMA-ES (2026)~\cite{robotnoise} & Sampling-time allocation based on sorting difficulty. & Motivates precision-aware evaluation; continuous robotics results are not prompt benchmarks.\\
\end{longtable}
\endgroup
